\documentclass[11pt, a4paper, openany]{scrartcl}
\usepackage{xcolor}
\usepackage[colorlinks=true,allcolors=blue,pdfborder={0 0 0},pdfstartview={FitV},breaklinks,linktocpage]{hyperref}
\usepackage{color,amsmath,amsfonts,amssymb}
\usepackage[utf8]{inputenc}
\usepackage{enumitem}%für mehrspaltige Aufzählungen
\usepackage[ngerman, english]{babel}
\usepackage{epsfig}
\usepackage{pstricks}
\usepackage{graphics}
\graphicspath{{Figures/}}
\usepackage{bbm}
%\usepackage{showlabels}
\usepackage[labelfont=bf]{caption}
%\usepackage{nonfloat}
\usepackage{booktabs}%für Tabellen
\usepackage{pgfplots} %für Plots
\pgfplotsset{
compat=1.5,
tick label style={font=\small},
every axis/.append style={line width=0.8pt, tick style={line width=0.7pt}}
} %für richtigen Abstand der labels
\usetikzlibrary{shapes,arrows,fit,calc,positioning,arrows}
\usepackage[
  textheight=25cm,
left=2.3cm,
right=2.3cm,
headheight=25pt,
  includehead,includefoot,
  heightrounded,
]{geometry}
\makeatletter
\def\input@path{{./Figures/}}
\usepackage{fancyhdr}
\pagestyle{fancy}
\fancyfoot{}
\fancyhead{}
% \renewcommand{\chaptermark}[1]
% {\markboth{{\chaptername\ \thechapter:\  #1}}{}}
% \renewcommand{\sectionmark}[1]
% {{\markright{ \thesection\ #1}}} 
%\renewcommand{\sectionmark}[1]{\markboth{#1}{}}

%\renewcommand\sectionmark[1]{\markright{#1}}
%\renewcommand{\thesection}{\arabic{section}.}
\fancyhead[L]{}
\fancyhead[R]{\thepage}
\fancyhead[C]{\nouppercase{\leftmark}}
%\fancyhead[C]{\nouppercase{\rightmark}}


%\usepackage{sectsty}
\usepackage{url}
%\usepackage{breakurl}
\usepackage{amsthm}
\usepackage{graphicx}
\usepackage{tikzpagenodes}

\definecolor{darkgreen}{rgb}{0.0,0.6,0.0}
\definecolor{shadedgreen}{rgb}{0.8,0.95,0.8}
\definecolor{shadedblue}{rgb}{0.85,0.85,1.0}
\newlength\figureheight 
\newlength\figurewidth


\usepackage[varg]{txfonts} %font/Schrift
\usepackage{csquotes}
\usepackage[minnames=3,maxnames=30,backend=biber,natbib=true,sorting=none,citestyle=numeric-comp,giveninits=true]{biblatex}
\addbibresource{Literature.bib}
%\setlength\bibitemsep{0.5em}
\setlength{\biblabelsep}{0.4em}

\DeclareFieldFormat{pages}{#1}
\DeclareFieldFormat{journaltitle}{#1}
\DeclareFieldFormat{eprint}{#1}
\DeclareFieldFormat[article]{title}{\it#1\isdot}

\DeclareBibliographyDriver{article}{
\printnames{author}, \printfield{title},
\href{http://dx.doi.org/\thefield{doi}}
{\printfield{journaltitle} 
\textbf{\printfield{volume}}, \printfield{pages} 
(\printfield{year})}.}

\DeclareBibliographyDriver{misc}{
\printnames{author}, \printfield{title},
\href{https://arxiv.org/abs/\thefield{eprint}}{
arXiv preprint, arXiv:\printfield{eprint}  (\printfield{year})}.}

\DeclareBibliographyDriver{inbook}{\hspace{-0.1cm}
\printnames{author},
\printfield{title},
in \printfield{booktitle}.
\printlist{publisher}
(\printfield{year}).
}


\renewbibmacro*{name:andothers}{% Based on name:andothers from biblatex.def
  \ifboolexpr{
    test {\ifnumequal{\value{listcount}}{\value{liststop}}}
    and
    test \ifmorenames
  }
    {\ifnumgreater{\value{liststop}}{1}
       {\finalandcomma}
       {}%
     \andothersdelim\bibstring[\emph]{andothers}}
    {}}

\setenumerate[1]{label=[\arabic*],ref=\arabic*}
\newcommand{\paper}[4]{\item #1, \,\textit{#2}, \,\href{#3}{#4}.\\[-1.4em]}


\newcommand{\bE}{\mathbf{E}}
\newcommand{\bA}{\mathbf{A}}
\newcommand{\eqt}{\,{=}\,}



\begin{document}

\pagenumbering{Roman}

\begin{titlepage}\pdfbookmark[0]{Title}{Title}
  \sffamily
  \begin{center}
{
\includegraphics[width=9cm]{logo.pdf}
\\[5em]
\Huge \bfseries Output summary of the CUED program}
\\[3em]\large
The CUED program is developed and maintained by:
\\[3em]
Computational Condensed Matter Theory Group
\\[0.5em]
Head: Ferdinand Evers
  \\[0.5em]
Institute of Theoretical Physics
  \\[0.5em]
University of Regensburg
  \\[0.5em]
Universitätsstraße 31
  \\[0.5em]
D\,-\,93053 Regensburg
  \\[0.5em]
Germany
\\[3em]
Version: Git commit GITCOMMITNUMBER, \today
\\[3em]
The calculation has been executed in the directory: 
\\[0.5em]
PH-DIRECTORY
\\[3em]
  \end{center}{\large
  With contributions from (in alphabetic order): 
  \\[1em]
  Jack Crewse
  \\[0.5em]
  Patrick Grössing
  \\[0.5em]
  Maximilian Nitsch
  \\[0.5em]
  Adrian Seith
  \\[0.5em]
  Jan Wilhelm
  }
\end{titlepage}







\pagenumbering{arabic}
\pagestyle{plain}




\pdfbookmark[0]{Contents}{Contents}
\tableofcontents


\pagestyle{fancy}

\section{Electric-field pulse}
The following electric driving field is employed in the simulation:
\begin{align}
    \bE(t)  = E(t)\,\hat{e}_\phi
    \,,\hspace{2em}
    E(t) = E_0\,    \sin\Big(2\pi f_0\,(1+f_\text{chirp} t) + \varphi\Big)\, e^{-t^2/(2\alpha)^2}\,.
    \label{e1}
\end{align}
\begin{figure}[b!]
\centering
\setlength\figureheight{8cm} 
\setlength\figurewidth{8cm}
\input{Efield.tikz}\hfill\input{Afield.tikz}
\caption{Left: Electric driving field~$E(t)$ from Eq.~\eqref{e1}, right: Gauge field~$A(t)$ from Eq.~\eqref{e2}.}
    \label{fig:Efield}
\end{figure}
The pulse is sketched in Fig.~\ref{fig:Efield}. 
%
The following parameters are used in the simulation:
\begin{itemize}
    \item PH-EFIELD-DIRECTION
    \item Pulse frequency: $f_0 = PH-FREQ$\,THz
    \item Chirp: $f_\text{chirp} = PH-CHIRP$\,THz
    \item Carrier-envelope phase: $\phi = PH-CEP$
    \item $\alpha = PH-ALPHA$\,fs, full width at half maximum (FWHM) = PH-FWHM\,fs
\end{itemize}
The gauge field~$\bA(t)\eqt A(t)\,\hat{e}_\phi$ follows from~$\dot\bA(t)\eqt -\bE(t)$ such that we have 
\begin{align}
    \dot A(t) = -E(t) \hspace{2em}\Rightarrow\hspace{2em}
    A(t) = -\int\limits_{-\infty}^t E(t')\,dt'\,. \label{e2}
\end{align}
The gauge field is sketched in Fig.~\ref{fig:Efield}. 


\section{Brillouin zone and \textit{k}-point grid}
REPLACEMENTBZ
\begin{figure}[b!]
\centering
\setlength\figureheight{16cm} 
\setlength\figurewidth{16cm}
% This file was created with tikzplotlib v0.10.1.
\begin{tikzpicture}

\definecolor{darkgray176}{RGB}{176,176,176}
\definecolor{gray}{RGB}{128,128,128}
\definecolor{green}{RGB}{0,128,0}

\begin{axis}[
height=\figureheight,	 scale only axis=true,
tick align=outside,
tick pos=left,
width=\figurewidth,
x grid style={darkgray176},
xlabel={\(\displaystyle k_x \; (\si{\per\As})\)},
xmin=-0.519674684499713, xmax=0.519674684499713,
xtick style={color=black},
y grid style={darkgray176},
ylabel={\(\displaystyle k_y \; (\si{\per\As})\)},
ymin=-0.141729459409013, ymax=0.141729459409013,
ytick style={color=black}
]
\addplot [thick, black, mark=*, mark size=1, mark options={solid}, only marks]
table {%
0 0
};
\addplot [thick, black]
table {%
0.472431531363375 0.094486306272675
-0.472431531363375 0.0944863062726751
-0.472431531363375 -0.094486306272675
0.472431531363375 -0.094486306272675
0.472431531363375 0.094486306272675
};
\addplot [thick, black, mark=*, mark size=1, mark options={solid}, only marks]
table {%
0.472431531363375 2.08166817117217e-17
0 0.0944863062726751
};
\addplot [thick, gray, mark=*, mark size=1, mark options={solid}, only marks]
table {%
-0.45353427010884 0.0472431531363375
-0.41573974759977 0.0472431531363375
-0.3779452250907 0.0472431531363375
-0.34015070258163 0.0472431531363375
-0.30235618007256 0.0472431531363375
-0.26456165756349 0.0472431531363375
-0.22676713505442 0.0472431531363375
-0.18897261254535 0.0472431531363375
-0.15117809003628 0.0472431531363375
-0.11338356752721 0.0472431531363375
-0.07558904501814 0.0472431531363375
-0.03779452250907 0.0472431531363375
0 0.0472431531363375
0.03779452250907 0.0472431531363375
0.0755890450181401 0.0472431531363375
0.11338356752721 0.0472431531363375
0.15117809003628 0.0472431531363375
0.18897261254535 0.0472431531363375
0.22676713505442 0.0472431531363375
0.26456165756349 0.0472431531363375
0.30235618007256 0.0472431531363375
0.34015070258163 0.0472431531363375
0.3779452250907 0.0472431531363375
0.41573974759977 0.0472431531363375
0.45353427010884 0.0472431531363375
-0.45353427010884 0.0472431531363375
};
\addplot [thick, gray, mark=*, mark size=1, mark options={solid}, only marks]
table {%
-0.45353427010884 -0.0472431531363375
-0.41573974759977 -0.0472431531363375
-0.3779452250907 -0.0472431531363375
-0.34015070258163 -0.0472431531363375
-0.30235618007256 -0.0472431531363375
-0.26456165756349 -0.0472431531363375
-0.22676713505442 -0.0472431531363375
-0.18897261254535 -0.0472431531363375
-0.15117809003628 -0.0472431531363375
-0.11338356752721 -0.0472431531363375
-0.07558904501814 -0.0472431531363375
-0.03779452250907 -0.0472431531363375
0 -0.0472431531363375
0.03779452250907 -0.0472431531363375
0.0755890450181401 -0.0472431531363375
0.11338356752721 -0.0472431531363375
0.15117809003628 -0.0472431531363375
0.18897261254535 -0.0472431531363375
0.22676713505442 -0.0472431531363375
0.26456165756349 -0.0472431531363375
0.30235618007256 -0.0472431531363375
0.34015070258163 -0.0472431531363375
0.3779452250907 -0.0472431531363375
0.41573974759977 -0.0472431531363375
0.45353427010884 -0.0472431531363375
-0.45353427010884 -0.0472431531363375
};
\addplot [thick, green]
table {%
0.383948684086951 0.118107882840844
-0.142340736398547 0.118107882840844
};
\addplot [thick, red]
table {%
-0.496053107931544 0.118107882840844
-0.142340736398547 0.118107882840844
};
\addplot [thick, black, mark=*, mark size=1, mark options={solid}, only marks]
table {%
-0.142340736398547 0.118107882840844
};
\draw (axis cs:0.01,0.01) node[
  scale=1,
  anchor=base west,
  text=black,
  rotate=0.0
]{$\Gamma$};
\draw (axis cs:0.472431531363375,2.08166817117217e-17) node[
  scale=1,
  anchor=base west,
  text=black,
  rotate=0.0
]{$\mr{X}$};
\draw (axis cs:0,0.0944863062726751) node[
  scale=1,
  anchor=base west,
  text=black,
  rotate=0.0
]{$\mr{Y}$};
\end{axis}

\end{tikzpicture}

\caption{Brillouin zone.}
    \label{fig:Efield}
\end{figure}



\section{Hamiltonian, bandstructure and dipoles}
REPLACEMENTHAMILTONIAN


\section{Time evolution of the density matrix}

\section{Time-dependent currents}

\section{Frequency-resolved emission spectrum}

\section{Summary of input parameters}

\section{Timings}


\section{References}
\begin{enumerate}[leftmargin=*]

\paper{J.~Wilhelm, P.~Grössing, A.~Seith, J.~Crewse, M.~Nitsch, L.~Weigl, C.~Schmid, and F.~Evers}{Semi\-con\-duc\-tor-Bloch Formalism: Derivation and Application to High-Harmonic Generation from Dirac Fermions}{TOBEFILLED}{ 
Phys.~Rev.~B~\,\textbf{x}, y (2021)}
\label{Wilhelm2021}

\end{enumerate}

\end{document}
